\begin{resumo}
Em aplicativos de transporte, a distância percorrida até o embarque (aproximação) é quilometragem sem passageiro, o \textit{km vazio}. Este trabalho avalia se a atribuição conjunta das solicitações já recebidas a veículos elegíveis, resolvida pelo algoritmo Húngaro em lotes de 60 segundos, reduz o km vazio e preserva o atendimento em comparação com o despacho pelo veículo mais próximo, imediato ou em lotes. A avaliação usa simulação contrafactual com entradas históricas de uma operação da Grande Vitória (ES) e, separadamente, 22 cenários sintéticos; as duas linhas não são somadas. Com a configuração congelada antes da execução, a atribuição conjunta reduziu o km vazio em relação ao despacho imediato em 5,96\% (15 dias de setembro de 2026) e 5,04\% (25 dias de outubro) e, em relação ao despacho em lotes pelo mais próximo, em 2,74\% e 1,40\%. O critério de atendimento foi preservado em setembro (16 atendimentos a mais), mas não em outubro, quando houve 85 atendimentos a menos que no despacho imediato (−0,36 ponto percentual do universo elegível); contra o despacho em lotes, foi preservado nos dois períodos. Nos 22 cenários sintéticos, que compartilham estrutura e não são confirmações independentes, a redução mediana foi de 8,4\%, com atendimento preservado em todos pelo critério sintético, mais tolerante que o histórico; essa evidência é exploratória e complementar. Os resultados são simulados, dependem de distância estimada e de um despacho de referência simplificado e não medem economia na operação real.
\end{resumo}

\begin{abstract}
In ride-hailing, the distance driven to reach the passenger (pickup) is mileage without a passenger, or \textit{empty kilometers}. This study assesses whether jointly assigning already-received requests to eligible vehicles, solved by the Hungarian algorithm in 60-second batches, reduces empty kilometers and preserves service compared with nearest-vehicle dispatching, either immediate or batched. The assessment uses counterfactual simulation with historical inputs from an operation in Greater Vitória (Brazil) and, separately, 22 synthetic scenarios; the two lines are never added together. With the configuration frozen before execution, joint assignment reduced empty kilometers relative to immediate dispatching by 5.96\% (15 days of September 2026) and 5.04\% (25 days of October), and relative to nearest-vehicle batching by 2.74\% and 1.40\%. The service criterion was preserved in September (16 more requests served) but not in October, when 85 fewer requests were served than under immediate dispatching (−0.36 percentage point of the eligible requests); against batching it was preserved in both periods. In the 22 synthetic scenarios, which share one structure and are not independent confirmations, the median reduction was 8.4\%, with service preserved in all of them under the synthetic criterion, which is more tolerant than the historical one; this evidence is exploratory and complementary. The results are simulated, rely on estimated distances and a simplified reference dispatcher, and do not measure savings in the real operation.
\end{abstract}
